\clearpage
\section*{Limitations}
\addcontentsline{toc}{section}{Limitations}
\label{sec:limitations}

\paragraph{Scope and external validity.}
Our conclusions concern one generator family, one coding-agent-style harness, and deterministic subsets of three long-form benchmarks.
The subset rule emphasizes long inputs and, when specified, long requested outputs rather than an independent measure of difficulty or of the need to revise.
The results therefore do not establish transfer to shorter tasks, other model families or runtimes, interactive human--AI writing, or end-to-end agents with retrieval.
The next test of adaptive process selection is a task set whose requests arrive during drafting or call for more drafting or review rather than replanning, so that rerunning a fixed cycle would choose the wrong process.

\paragraph{Comparison design.}
The alternative writing systems are reimplementations under the shared setup of Section~\ref{sec:implementation}, not executions of the published systems.
\condTaskPlanning\ also delegates its task nodes to subagents, so delegation is present on both sides of that comparison.
Its prompt forbids rewriting already-composed text, so \condTaskPlanning\ has no revision pass over the assembled draft.
The \condExploratoryTwo\ adaptation runs only on ten consensus-writing tasks, and running it on the three benchmarks is the next comparison.\footnote{See Appendix~\ref{app:habermas} for the consensus-writing results.}
The compute-matched analysis selects pairs with similar realized token use after generation, and the post-hoc \condSelfRefine\ baseline, budgeted before generation on unscored prompts, received \SelfRefineTokenExcess\ more output-plus-reasoning tokens per document than our system (\SelfRefineOutputTokens\ against \SelfRefineAfourOutputTokens).

\paragraph{Post-hoc conditions.}
\condStagedRoles\ was added post-hoc, uses one generation run and one judge family, and differs from our system in the \texttt{Monitor}, goals, shared state, and delegation together, so it isolates the role-instruction content and no single component.
\condPlainAgent\ was also added post hoc on one generation run; it shares our harness and generator but no skill content.
\condExploratoryOne\ on the three benchmarks is our adaptation of CogWriter, run once post hoc: it completed \CogWriterStyleCompleted\ of \CogWriterStyleAttempted\ prompts, mostly failing its final-product gate for lack of a draft file, only \CogWriterStyleRetried\ of its \CogWriterStyleFirstPassFailed\ first-pass failures were retried before a cost guard stopped retries, and it used more compute than our system (\CogWriterStyleSpawns\ against \AfourSpawns\ subagent spawns per run).
\condSelfRefine\ and \condLoopOnly\ also ran once post hoc, completing \SelfRefineCompleted\ and \LoopOnlyCompleted\ of \CogWriterStyleAttempted\ prompts; \condLoopOnly\ is built from the skill texts that \condNoGoals, \condFixedOrder, and \condSingleContext\ executed in generation run 1.

\paragraph{Output language.}
On WritingBench, some English assignments are answered mainly in Chinese in every condition, from \AtwoDriftCount\ of \AtwoDriftEnglish\ for \condStaged\ to \AfiveDriftCount\ of \AfiveDriftEnglish\ for \condNoGoals\ in generation run 1, with \AfourDriftCount\ of \AfourDriftEnglish\ for our system.
The drift concentrates on prompts with Chinese topics or instructions and varies by condition.\footnote{See Appendix~\ref{app:sensitivity-analyses} for every condition.}

\paragraph{Missing outputs.}
Headline win rates exclude prompts without a completed output; counting every attempted prompt in the denominator, with ties and missing outputs as non-wins, our system wins \UncondFullSinglePassWinShare\ (\condSinglePass), \UncondFullStagedWinShare\ (\condStaged), and \UncondFullTaskPlanningWinShare\ (\condTaskPlanning) of prompts, and assigning every missing pair against our system leaves win rates over decided pairs of at least \WorstCaseFullSinglePassRate\ (\condSinglePass), \WorstCaseFullStagedRate\ (\condStaged), and \WorstCaseFullTaskPlanningRate\ (\condTaskPlanning).
Some systems fail more prompts than others, so this accounting bounds the effect of the excluded pairs.

\paragraph{Evaluator dependence.}
The primary evaluator belongs to the same model family as the generator, and related generator--judge models can exhibit preference leakage or self-preference~\citep{li2026preference,panickssery2024selfpreference}.
The cross-family judge rescored generation run 1 for the five contrasts between our system and the pre-specified alternatives and ablations, and for no post-hoc condition; summing W/L/T counts over the three benchmarks favors our system over \condStaged\ (\CrossFamilyFullStagedWinRate, $p=\CrossFamilyFullStagedP$) and leaves \condFixedOrder\ near parity (\CrossFamilyFullFixedOrderWinRate, $p=\CrossFamilyFullFixedOrderP$), as under the primary judge.
The judge produces many ties and reverses 5 of 15 benchmark-level comparisons, and its agreement with the primary judge is $\kappa=\JudgeAgreementKappa$ over all \JudgeAgreementPairs\ prompt-level outcomes and $\kappa=\JudgeAgreementDecidedKappa$ over the \JudgeAgreementDecidedPairs\ pairs both judges decide; extending it to all three generation runs is the next step.
Model-judge preference is not human preference; a blinded human comparison of same-prompt outputs, stratified by benchmark, contrast, and prompt length, would measure judge agreement with human readers.

\paragraph{Component attribution.}
The ablations constrain explicit goals and adaptive ordering to modest contributions without establishing equivalence.
Neither ablation differs from the full system in any benchmark--run test after Holm correction, each such test reaches 80\% power~\citep{card2020little} at a true preference of about \PowerCellDetectablePreference, and the post-hoc prompt-level estimate compares each ablation only with the full system.\footnote{See Appendix~\ref{app:prompts} for the power calculation.}
\condFixedOrder\ enforces its cycle through the \texttt{Monitor}'s instructions so that it shares our system's agent runtime and differs only in the selection rule, and \AsixCycleCompliantRate\ of its paths are cycle-compliant.\footnote{See Appendix~\ref{app:process-behavior} for the trace classification.}
The combined \condLoopOnly\ ablation is post hoc and from one generation run, and its parity with our system comes without an equivalence test; it does not separate the \texttt{Monitor}, the workspace files, in-place editing, and the role guidance, which are the next ablations, together with a runner-enforced cycle.

\paragraph{Execution variants.}
\condSingleWriter\ changes delegation, role context, decision granularity, and output length together; \condSingleContext\ is narrower but was added post-hoc, uses one exploratory generation run, and differs in execution interface and transport.
The current evidence therefore does not attribute the remaining system-level advantage to any single shared execution feature, and delegation itself varies: the \texttt{Monitor} performed at least one role itself when native delegation was unavailable in \AfourFallbackRuns\ of \AfourRunCount\ completed runs of our system, and in \AfiveFallbackRuns\ and \AsixFallbackRuns\ runs of \condNoGoals\ and \condFixedOrder, although dropping them barely changes the comparison with the alternatives.\footnote{See Appendix~\ref{app:process-behavior} for the fallback analysis.}

\paragraph{Descriptive execution analyses.}
The primary evaluator favors longer outputs overall, choosing the longer output in \LengthRuleAcrossRate\ of decided pairs, and the length- and compute-matched analyses are descriptive because both variables result from execution.\footnote{See Appendix~\ref{app:sensitivity-analyses} for the rate per benchmark.}
Our system also costs more time: it averages \AfourMeanSecondsMean\ wall-clock seconds per document against \AoneMeanSecondsMean\ for \condSinglePass, while \condNoGoals\ (\AfiveMeanSecondsMean\ seconds) and \condSingleContext\ (\SingleContextMeanSeconds\ against \AfourMeanSeconds\ on its generation run) stay near parity with it at lower cost.\footnote{See Appendices~\ref{app:process-behavior} and~\ref{app:sensitivity-analyses} for the time and token accounting.}
Process traces record state transitions and agent actions; they characterize how \condAgenticCogWriter\ behaves on these tasks but do not measure human cognition or validate cognitive process theory as a model of language-agent internals.
